\documentclass[12pt,a4paper]{article}
\usepackage[utf8]{inputenc}
\usepackage[T1]{fontenc}
\usepackage[french]{babel}
\usepackage[a4paper,margin=2.2cm]{geometry}
\usepackage{xcolor}
\usepackage{listings}

\lstdefinestyle{pythonstyle}{
    language=Python,
    basicstyle=\ttfamily\small,
    breaklines=true,
    showstringspaces=false,
    upquote=true
}

\lstdefinestyle{terminalstyle}{
    backgroundcolor=\color{black},
    basicstyle=\ttfamily\small\color{white},
    frame=none,
    breaklines=true,
    showstringspaces=false,
    xleftmargin=10pt,
    xrightmargin=10pt,
    aboveskip=0pt
}

\begin{document}

\section*{Exercice 1 }
\begin{lstlisting}[style=pythonstyle]
notes = [12, 15, 9, 18, 7, 14]
print(notes)
print(len(notes))
if len(notes) > 5:
    print("Plus de 5 elements")
else:
    print("5 elements ou moins")
for n in notes:
    print(n)
\end{lstlisting}
\begin{lstlisting}[style=terminalstyle]
[12, 15, 9, 18, 7, 14]
6
Plus de 5 elements
12
15
9
18
7
14
\end{lstlisting}

\section*{Exercice 2 }
\begin{lstlisting}[style=pythonstyle]
liste = list(range(1, 11))
print(liste)
for n in liste:
    if n % 2 == 0:
        print(n)
c = 0
for n in liste:
    if n >= 6:
        c = c + 1
print(c)
\end{lstlisting}
\begin{lstlisting}[style=terminalstyle]
[1, 2, 3, 4, 5, 6, 7, 8, 9, 10]
2
4
6
8
10
5
\end{lstlisting}
\clearpage
\section*{Exercice 3}
\begin{lstlisting}[style=pythonstyle]
depenses = [25, 40, 18, 32, 15]
print(sum(depenses))
print(sum(depenses) / len(depenses))
if sum(depenses) > 100:
    print("Somme > 100")
c = 0
for d in depenses:
    if d > 20:
        c = c + 1
print(c)
\end{lstlisting}
\begin{lstlisting}[style=terminalstyle]
130
26.0
Somme > 100
3
\end{lstlisting}

\section*{Exercice 4}
\begin{lstlisting}[style=pythonstyle]
temperatures = [18, 22, 17, 25, 19, 30, 16]
print(min(temperatures))
print(max(temperatures))
print(max(temperatures) - min(temperatures))
for t in temperatures:
    if t > 20:
        print(t)
\end{lstlisting}
\begin{lstlisting}[style=terminalstyle]
16
30
14
22
25
30
\end{lstlisting}

\section*{Exercice 5}
\begin{lstlisting}[style=pythonstyle]
fruits = ["pomme", "banane", "orange", "fraise", "mangue"]
print("banane" in fruits)
print("kiwi" in fruits)
f = input()
if f in fruits:
    print("Present")
else:
    print("Absent")
\end{lstlisting}
\begin{terminalstyle}
True
False
fraise (saisie utilisateur)
Present
\end{terminalstyle}
\clearpage
\section*{Exercice 6}
\begin{lstlisting}[style=pythonstyle]
liste1 = [4, 8, 15, 16, 23, 42]
liste2 = liste1.copy()
print(liste1)
print(liste2)
liste2[0] = 0
print(liste1)
print(liste2)
for x in liste2:
    print(x)
\end{lstlisting}
\begin{lstlisting}[style=terminalstyle]
[4, 8, 15, 16, 23, 42]
[4, 8, 15, 16, 23, 42]
[4, 8, 15, 16, 23, 42]
[0, 8, 15, 16, 23, 42]
0
...
\end{lstlisting}

\section*{Exercice 7}
\begin{lstlisting}[style=pythonstyle]
valeurs = [7, 12, 5, 18, 3, 12, 9]
print(len(valeurs))
print(sum(valeurs))
print(min(valeurs))
print(max(valeurs))
print(12 in valeurs)
c = 0
for v in valeurs:
    if v > 10:
        c = c + 1
print(c)
\end{lstlisting}
\begin{lstlisting}[style=terminalstyle]
7
66
3
18
True
3
\end{lstlisting}
\clearpage
\section*{Exercice 8}
\begin{lstlisting}[style=pythonstyle]
nombres = list(range(1, 21))
print(nombres)
print(len(nombres))
print(sum(nombres))
p = []
for n in nombres:
    if n % 2 == 0:
        p.append(n)
print(p)
print(min(p))
print(max(p))
\end{lstlisting}
\begin{lstlisting}[style=terminalstyle]
[1, ..., 20]
20
210
[2, 4, ..., 20]
2
20
\end{lstlisting}

\section*{Grand Exercice 1}
\begin{lstlisting}[style=pythonstyle]
notes = [11, 14, 9, 16, 7, 18, 13, 10, 5, 19]
print(len(notes))
print(sum(notes))
print(sum(notes)/len(notes))
print(min(notes))
print(max(notes))
print(20 in notes)
c = 0
for n in notes:
    if n >= 10:
        c = c + 1
print(c)
cp = notes.copy()
for i in range(len(cp)):
    if cp[i] < 10:
        cp[i] = 10
print(notes)
print(cp)
\end{lstlisting}
\begin{lstlisting}[style=terminalstyle]
10
122
12.2
5
19
False
7
[11, 14, 9, 16, 7, 18, 13, 10, 5, 19]
[11, 14, 10, 16, 10, 18, 13, 10, 10, 19]
\end{lstlisting}

\section*{Grand Exercice 2}
\begin{lstlisting}[style=pythonstyle]
stock = [12, 5, 8, 0, 17, 3, 0, 9]
print(len(stock))
print(sum(stock))
print(min(stock))
print(max(stock))
r = 0
for q in stock:
    if q == 0:
        r = r + 1
print(r)
print(17 in stock)
cp = stock.copy()
for i in range(len(cp)):
    if cp[i] < 10:
        cp[i] = cp[i] + 5
print(stock)
print(cp)
d = []
for q in stock:
    if q > 0:
        d.append(q)
print(d)
\end{lstlisting}
\begin{lstlisting}[style=terminalstyle]
8
54
0
17
2
True
[12, 5, 8, 0, 17, 3, 0, 9]
[12, 10, 13, 5, 17, 8, 5, 14]
[12, 5, 8, 17, 3, 9]
\end{lstlisting}

\end{document}